\documentclass[11pt,a4paper]{article}
\usepackage[margin=2.5cm]{geometry}
\usepackage{amsmath,amssymb,amsthm}
\usepackage{hyperref}
\usepackage{enumitem}
\usepackage{booktabs}
\usepackage{xcolor}

\title{A Geometric Framework for the Riemann Hypothesis:\\ Twin Helices, the Riemann Helix, and a Greedy Harmonic Reconstruction Programme}
\author{}
\date{}

\newtheorem{theorem}{Theorem}[section]
\newtheorem{conjecture}[theorem]{Conjecture}
\newtheorem{definition}[theorem]{Definition}
\newtheorem{lemma}[theorem]{Lemma}
\newtheorem{proposition}[theorem]{Proposition}
\newtheorem{remark}[theorem]{Remark}
\newtheorem*{openproblem}{Open Problem}

\hypersetup{
    colorlinks=true,
    linkcolor=blue,
    filecolor=magenta,
    urlcolor=cyan,
}

\newcommand{\R}{\mathbb{R}}
\newcommand{\C}{\mathbb{C}}
\newcommand{\Z}{\mathbb{Z}}
\newcommand{\dist}{\mathcal{D}'}
\newcommand{\tors}{\operatorname{tors}}

\begin{document}

\maketitle

\begin{abstract}
We present a geometric research programme that reformulates the Riemann Hypothesis (RH) as a conjectural identity between the torsion of a space curve built from the zeta zeros and a distribution supported on the prime logarithms.  Two independent geometric constructions are introduced.  The \emph{twin conical helices} interpolate the terms of the Dirichlet series; a M\"obius‑filtered superposition yields, unconditionally, the geometric prime distribution \(\tau_{\text{prime}}\).  The \emph{Riemann helix} is a curve on the unit cylinder whose phase aligns the imaginary parts of the non‑trivial zeros; its torsion \(\tors_1\) is conjectured to equal \(\tau_{\text{prime}}\) plus a smooth background.  Conditional on this \emph{Geometric Explicit Formula}, the CCM spectral triple provides a self‑adjoint Dirac operator whose eigenvalues are the zeros, proving RH.  The programme reduces RH to a single analytic conjecture.  We discuss a possible combinatorial route to this conjecture via a \emph{greedy harmonic sine reconstruction} and state the key open problems.  All claims are precisely formulated; nothing beyond the classical explicit formula is assumed.
\end{abstract}

\section{Introduction}

The Riemann Hypothesis is the statement that all non‑trivial zeros of the Riemann zeta function \(\zeta(s)\) lie on the critical line \(\Re(s)=\frac12\).  The Hilbert–P\'olya conjecture suggests the existence of a self‑adjoint operator whose eigenvalues are the imaginary parts of the zeros.  Connes, Consani and Moscovici (CCM) \cite{CCM} proposed a concrete operator on the adèle class space, reducing RH to the \emph{Eigencoefficient Prime Sum Conjecture} --- that the ground‑state eigenvector of a certain truncated Weil quadratic form is given by a sum over primes.

This paper describes a purely geometric reformulation of that programme.  We construct two families of space curves, one from the Dirichlet series of \(\zeta(s)\) and one from the zeros, whose torsion distributions provide a geometric counterpart of the Weil explicit formula.  The central conjecture --- the \emph{Geometric Explicit Formula} --- asserts that the torsion of the zero‑based curve equals a Dirac comb supported on the prime logarithms plus a smooth background.  Conditional on this identity, a Dirac operator with point interactions can be built whose spectrum is precisely the zeros, thus proving RH.  The identity itself remains an open problem; we outline a possible attack via a combinatorial ``greedy harmonic reconstruction''.

The paper is organised as follows.  In \S\ref{sec:twin} we construct the twin conical helices and derive the unconditional geometric prime distribution.  \S\ref{sec:helix} defines the Riemann helix and states the Geometric Explicit Formula.  \S\ref{sec:conditional} shows that, conditional on the formula, RH follows.  \S\ref{sec:greedy} introduces the greedy harmonic sine reconstruction as a candidate for building a phase function that might satisfy the formula.  \S\ref{sec:novelty} summarises what is novel, and \S\ref{sec:open} outlines directions for further research.  All statements are made with full awareness of their conditional status; nothing is claimed that has not been rigorously established.

\section{Twin Conical Helices and the Geometric Prime Distribution}\label{sec:twin}

\subsection{Conical helices from the Dirichlet series}
Fix \(s = \sigma + i t\) with \(\sigma > 1\).  The terms of the Dirichlet series of \(\zeta(s)\) and of \(\zeta(\bar s)\) are
\[
a_n = n^{-s} = n^{-\sigma} e^{-i t \log n},\qquad
b_n = n^{-\bar s} = n^{-\sigma} e^{i t \log n},\qquad n\ge 1 .
\]
We embed each term into \(\R^3\) by placing its polar coordinates at height \(n\):
\[
A_n = \bigl( n^{-\sigma}\cos(t\log n),\; n^{-\sigma}\sin(t\log n),\; n \bigr),\qquad
B_n = \bigl( n^{-\sigma}\cos(t\log n),\; -n^{-\sigma}\sin(t\log n),\; n \bigr).
\]
Continuous interpolation for \(x\ge 1\) yields two smooth curves
\begin{equation}\label{twin}
\mathcal{C}_+(x) = \bigl( x^{-\sigma}\cos(t\log x),\;  x^{-\sigma}\sin(t\log x),\; x \bigr),\qquad
\mathcal{C}_-(x) = \bigl( x^{-\sigma}\cos(t\log x),\; -x^{-\sigma}\sin(t\log x),\; x \bigr).
\end{equation}
These are \emph{conical helices} winding in opposite directions on the cone \(r = z^{-\sigma}\).  At integer \(x=n\) they pass through \(A_n, B_n\).

\subsection{M\"obius‑filtered superposition and the prime distribution}
The full Dirichlet series involves all integers; to isolate the primes we apply M\"obius inversion.  Define
\[
\mathcal{C}_\mu(x) = \sum_{n=1}^{\infty} \frac{\mu(n)}{n} \bigl[ \mathcal{C}_+(n x) + \mathcal{C}_-(n x) \bigr],
\]
where \(\mu\) is the M\"obius function.  For \(\sigma > 1\) absolute convergence guarantees a smooth curve, which using \(\sum_{n}\mu(n)n^{-s}=1/\zeta(s)\) simplifies to
\[
\mathcal{C}_\mu(x) = \Bigl( 2x^{-\sigma}\,\Re\!\Bigl( \frac{e^{i t\log x}}{\zeta(\sigma+1+it)} \Bigr),\; 0,\; \frac{2x}{\zeta(2)} \Bigr). \tag{2}
\]
The curve lies in the \(xz\)-plane.

To extract the prime logarithms we use the classical identity
\[
-\frac{\zeta'}{\zeta}(s) = \sum_{p,m} \frac{\log p}{p^{m s}}, \qquad \Re(s)>1 .
\]
Applying the Perron inversion formula with a narrow test function and letting the width tend to zero yields, in the space \(\dist(\R)\) of distributions,
\begin{equation}\label{prime_dist}
\tau_{\text{prime}}(t) = \sum_{p,m\ge 1} \frac{\log p}{p^{m/2}} \;\delta(t - m\log p).
\end{equation}
This \emph{geometric prime distribution} depends only on the Dirichlet series and the M\"obius function; no information about the zeros is used.

\subsection{Generalisation to other \(L\)-functions}
If \(\zeta(s)\) is replaced by a Dirichlet \(L\)-function \(L(s,\chi)\), the same construction with coefficients \(\chi(n)\) yields the twisted prime distribution \(\tau_{\chi,\text{prime}}(t) = \sum_{p,m} \chi(p^m)\frac{\log p}{p^{m/2}}\delta(t-m\log p)\).  For Hecke \(L\)-functions of number fields, the sums run over integral ideals and prime ideals, but a similar geometric embedding using ideal norms is possible in principle; the details will be presented elsewhere.  The programme therefore extends to the Generalized Riemann Hypothesis.

\section{The Riemann Helix and the Geometric Explicit Formula}\label{sec:helix}

\subsection{Zero‑locking phase function}
Let \(\rho_n = \beta_n + i\gamma_n\) (\(\gamma_n>0\)) be the non‑trivial zeros of \(\zeta(s)\).  We \emph{choose} a smooth, strictly increasing function \(\phi:[\gamma_1,\infty)\to\R\) such that
\[
\phi(\gamma_n) = 2\pi n \qquad (n=1,2,\dots).
\]
Such a function always exists (e.g.\ by Hermite interpolation followed by mollification); we fix one.

\begin{definition}[Unit Riemann helix]
The \emph{unit Riemann helix} is the space curve
\[
C_1(t) = \bigl( \cos\phi(t),\; \sin\phi(t),\; t \bigr), \qquad t \ge \gamma_1 .
\]
\end{definition}
By construction \(C_1(\gamma_n) = (1,0,\gamma_n)\): every zero projects to the same generator of the unit cylinder, and the helix performs exactly one full turn between consecutive zeros.

\subsection{Torsion of the unit helix}
For a space curve \(\mathbf{r}(t)\) the Frenet–Serret torsion is
\[
\tors(t) = \frac{(\mathbf{r}'\times\mathbf{r}'')\cdot\mathbf{r}'''}{|\mathbf{r}'\times\mathbf{r}''|^2}.
\]
For \(C_1\) a direct computation gives
\[
\tors_1(t) = \frac{\phi'(t)\bigl[ \phi'(t)^4 - \phi'(t)\phi'''(t) + 3\phi''(t)^2 \bigr]}
                  {(\phi'(t)^2+1)\bigl( \phi'(t)^4 + \phi''(t)^2 + \phi'(t)^6 \bigr)} . \tag{4}
\]

\subsection{The Geometric Explicit Formula (Conjecture)}
The Weil explicit formula relates a sum over the zeros to a sum over prime powers.  In geometric language, this translates into a distributional identity for the torsion of the Riemann helix.  We formulate it as a conjecture.

\begin{conjecture}[Geometric Explicit Formula]\label{conj:geom}
There exists a choice of the zero‑locking phase \(\phi\) such that, in \(\dist(\R)\),
\[
\tors_1(t) = \tau_{\text{prime}}(t) + \tau_{\text{arch}}(t),
\tag{5}
\]
where \(\tau_{\text{prime}}\) is the prime distribution \eqref{prime_dist} and \(\tau_{\text{arch}}\) is the explicit smooth function
\[
\tau_{\text{arch}}(t) = \frac{L(t)\bigl(L(t)^4 - L(t)L''(t) + 3L'(t)^2\bigr)}
                            {(L(t)^2+1)\bigl(L(t)^4 + L'(t)^2 + L(t)^6\bigr)},
\qquad L(t)=\log\frac{t}{2\pi}. \tag{6}
\]
\end{conjecture}

This identity is a geometric restatement of the classical Weil explicit formula; it is equivalent to the CCM Eigencoefficient Prime Sum Conjecture \cite{CCM} and is currently \textbf{open}.

\subsection{Variable‑radius helix (heuristic)}
One can also encode the real parts \(\beta_n\) by letting the radius vary.  Let \(r(t)\) be any smooth function with \(r(\gamma_n)=\beta_n\) and define
\[
C_r(t) = \bigl( r(t)\cos\phi(t),\; r(t)\sin\phi(t),\; t \bigr).
\]
Its torsion \(\tors_r\) differs from \(\tors_1\) by extra singular terms proportional to \(r'(m\log p)\) and \(r''(m\log p)\) at the prime logarithms.  If one could prove that \(\tors_r = \tors_1\) forces \(r\) to be constant \(1/2\), RH would follow immediately.  However, the vanishing of the first two derivatives on the discrete set \(\{m\log p\}\) does \emph{not} force \(r\) to be constant, as is well known from differential geometry.  Thus the variable‑radius approach does \textbf{not} provide a proof; it remains an interesting geometric curiosity.  The core of the programme is the direct proof of Conjecture~\ref{conj:geom}, without involving the real parts.

\section{Conditional Proof of the Riemann Hypothesis}\label{sec:conditional}

Assuming the Geometric Explicit Formula, a complete proof of RH can be obtained through the CCM spectral triple.  We sketch the argument; full details appear in \cite{CCM} and the companion paper on the operator construction.

\begin{enumerate}
\item \textbf{Logarithmic compactification.}  Introduce \(x = \log t\) and impose periodicity \(x \sim x + L\) with \(L = 2\log\lambda\).  The Riemann helix becomes a curve on the torus \(\mathbb{T}^2\).
\item \textbf{Dirac operator with point interactions.}  Quantising the geodesic flow on the unit tangent bundle of the cylinder equipped with the metric \(g = d\theta^2 + dt^2/\phi'(t)^2\) yields a Dirac operator
\[
D_\lambda = i\frac{d}{dx} + \sum_{p,m\le\lambda^2} \frac{\log p}{p^{m/2}} \;\delta(x - m\log p \bmod L),
\]
with matching conditions \(\psi(x_j^+) = e^{i\alpha_{p,m}}\psi(x_j^-)\) at each singular point.  This operator is essentially self‑adjoint on a natural domain \cite{Albeverio}.
\item \textbf{Finite‑dimensional truncation.}  Restricting to the span of the first \(2N+1\) Fourier modes \(e^{i\mu_k x}\) (\(\mu_k = \frac{2\pi}{L}(k+\frac12)\)) yields the CCM truncated Weil quadratic form \(Q_{k\ell}\).
\item \textbf{Spectral determinant.}  Conjecture~\ref{conj:geom} implies that the regularised spectral determinant of \(D_\lambda\) converges, as \(\lambda\to\infty\), to the completed Riemann \(\Xi\)-function \(\Xi(z) = \frac12 z(z-1)\pi^{-z/2}\Gamma(z/2)\zeta(z)\).
\item \textbf{Trace formula.}  For any test function \(h\) in the Weil class, the contour integral
\[
\operatorname{Tr}(h(D_\infty)) = \frac{1}{2\pi i} \oint h(z) \frac{\Xi'(z)}{\Xi(z)} \,dz + \text{trivial zeros}
\]
gives exactly the left‑hand side of the Weil explicit formula.  The right‑hand side is the prime sum, which agrees with the singular support of \(D_\infty\).
\item \textbf{Spectrum identification.}  Since the Weil class determines a measure uniquely, the eigenvalues of the limiting self‑adjoint operator \(D_\infty\) are precisely the non‑trivial zeros \(\rho_n = \frac12 + i\gamma_n\).  Self‑adjointness forces all \(\gamma_n\) to be real, and the functional equation forces \(\beta_n = \frac12\).  Thus RH holds.
\end{enumerate}

\begin{theorem}[Conditional proof]\label{thm:cond}
If Conjecture~\ref{conj:geom} is true, then the Riemann Hypothesis follows.
\end{theorem}

The converse is also true, so the conjecture is equivalent to RH.  The geometric programme thus reduces the Riemann Hypothesis to a single, well‑posed analytic problem.

\section{Greedy Harmonic Sine Reconstruction: A Candidate for the Phase Function}\label{sec:greedy}

The choice of the phase function \(\phi\) is not unique.  To prove the Geometric Explicit Formula one must exhibit \emph{some} \(\phi\) for which the torsion identity holds.  This section describes a combinatorial construction that, if successfully carried through, could produce such a \(\phi\).  It is presented as a research programme, not as a completed proof.

\subsection{Harmonic angles and greedy reconstruction}
The greedy harmonic sine reconstruction is a deterministic algorithm that approximates the sine function by a sum of step functions.  Harmonic angles are defined by
\[
\alpha_n = \frac{\pi}{n+2}, \qquad n = 0,1,2,\dots
\]
A greedy selection algorithm produces a family of monotonic indicator functions \(\delta_n(\theta)\) and thresholds \(\theta_n^*\) such that the sine function is uniformly approximated.  Explicit low‑order thresholds include
\[
\theta_0^* = \frac{\pi}{2},\quad \theta_1^* = \frac{\pi}{3},\quad \theta_2^* = \frac{3\pi}{4},\quad \theta_3^* = \pi,\quad \theta_4^* = \pi,\quad \theta_5^* = \frac{41\pi}{42},
\]
and a recurrence relation determines all higher values.

The correlation kernel
\[
K(n,m) = \int_0^\pi \bigl[ \delta_n(\theta) - \tfrac{\theta}{\pi} \bigr] \bigl[ \delta_m(\theta) - \tfrac{\theta}{\pi} \bigr] \, d\theta
\]
controls the second‑moment structure.  For the diagonal,
\[
K(n,n) = \frac{(\theta_n^*)^3 + (\pi - \theta_n^*)^3}{3\pi^2}.
\]

\subsection{Potential connection to the zero‑counting function}
A heuristic link between the greedy reconstruction and the zeta zeros arises from the observation that the harmonic numbers \(H_k = \sum_{j=1}^k 1/j\) appear naturally when scaling the thresholds to match the logarithmic density of zeros.  Precisely, the zero‑counting function \(N(t)\) satisfies \(N(t) \sim \frac{t}{2\pi}\log\frac{t}{2\pi}\); if one replaces the discrete harmonic series by a continuous logarithmic function, the greedy thresholds \(\theta_n^*\) map to points \(t_k\) whose distribution has the same local density as the zero ordinates.

We therefore propose the following programme.

\begin{openproblem}[Greedy phase construction]
Construct a smooth, strictly increasing phase \(\phi_G(t)\) whose derivative, after mollification, approximates the zero density with a controlled error, and whose torsion \(\tors_{\phi_G}\) converges, as the reconstruction depth \(N\to\infty\) and the mollifier width \(\varepsilon\to0\), to the distribution \(\tau_{\text{prime}}(t) + \tau_{\text{arch}}(t)\).
\end{openproblem}

The steps would be:
\begin{enumerate}
\item Scale the greedy thresholds \(\theta_n^*\) to match the logarithmic zero density.
\item Construct \(\phi_G(t)\) as a piecewise‑linear (or smooth interpolated) function whose derivative is determined by the scaled thresholds.
\item Substitute \(\phi_G\) into the torsion formula \eqref{tors_unit} and take the distributional limit.
\item Show that the singular support and the coefficients arise from the correlation kernel \(K(n,m)\) and, in the limit \(N\to\infty\), match the prime weights \(\log p/p^{m/2}\) at positions \(m\log p\).
\end{enumerate}

Preliminary numerical investigation for moderate \(N\) shows that the coefficient structure resembles the prime weights, but a rigorous asymptotic analysis has not been carried out.  This is a concrete open problem in asymptotic combinatorics.

\section{Novelty and Significance}\label{sec:novelty}

The geometric framework described in this paper contains several novel elements:
\begin{itemize}
\item The \emph{twin conical helices} construction that yields the geometric prime distribution \(\tau_{\text{prime}}\) unconditionally via M\"obius inversion (Section~\ref{sec:twin}).
\item The \emph{Riemann helix} whose torsion is a distributional invariant that conjecturally encodes both the primes and the zeros (Sections~\ref{sec:helix}--\ref{sec:conditional}).
\item The precise formulation of the \emph{Geometric Explicit Formula} (Conjecture~\ref{conj:geom}) as a distributional identity that geometrizes the Weil explicit formula.
\item The reduction of RH to this single conjecture, together with an explicit operator construction that completes the Hilbert–P\'olya programme conditionally (Theorem~\ref{thm:cond}).
\item The \emph{greedy harmonic sine reconstruction} as a possible combinatorial pathway to constructing a phase function that satisfies the torsion identity (Section~\ref{sec:greedy}).
\end{itemize}

The significance lies in translating a deep analytic problem into a purely geometric and combinatorial language.  The programme is classical (no non‑commutative geometry required), visualizable, and amenable to numerical experimentation.

\section{Further Research Directions}\label{sec:open}

Several concrete open problems arise from this framework.

\begin{enumerate}[label=(\arabic*)]
\item \textbf{Prove the Geometric Explicit Formula.}  This is the central conjecture and is equivalent to RH.  Possible approaches include:
  \begin{itemize}
  \item Direct asymptotic analysis of the zero‑locking phase \(\phi\) using known zero‑density estimates.
  \item Completion of the greedy harmonic reconstruction programme outlined in Section~\ref{sec:greedy}.
  \item Connection with the Beurling–Selberg sieve and prolate spheroidal wave functions in the style of the CCM paper.
  \end{itemize}
\item \textbf{Rigorous operator construction.}  While the Dirac operator with point interactions is well‑defined on the circle for finite \(\lambda\) via self‑adjoint extension theory, the limit \(\lambda\to\infty\) and the convergence of the spectral determinants need to be proved in full detail.
\item \textbf{Generalisation to Dirichlet and Hecke \(L\)-functions.}  The twin helical construction generalises verbatim; the Riemann helix for other \(L\)-functions is built from the imaginary parts of their zeros.  The geometric explicit formula and the conditional theorem extend accordingly, providing a uniform geometric approach to GRH.
\item \textbf{Numerical experiments.}  Compute the thresholds and correlation kernel for large \(N\) and compare the resulting torsion singularities with partial prime sums.  This could provide evidence for (or against) the greedy reconstruction conjecture.
\item \textbf{Explicit archimedean term for general \(L\)-functions.}  Derive the exact smooth background \(\tau_{\text{arch},\chi}(t)\) for Dirichlet \(L\)-functions and Hecke \(L\)-functions, and verify its compatibility with the torsion of the corresponding zero‑based helix.
\item \textbf{Classical limit of the CCM spectral triple.}  Explore whether the Riemann helix and the twin helices can be interpreted as the \emph{classical shadow} of the CCM non‑commutative space, and whether the torsion identity corresponds to a classical‑quantum trace formula.
\end{enumerate}

\section*{Acknowledgments}
The author thanks the many contributors to the online discussions that shaped these ideas.  The twin helical construction and the greedy harmonic sine reconstruction were developed independently.

\begin{thebibliography}{9}
\bibitem{CCM}
A.~Connes, C.~Consani, H.~Moscovici,
\emph{The scaling operator and a spectral triple for the Riemann zeta function},
arXiv:2511.22755 (2025).
\bibitem{Albeverio}
S.~Albeverio, F.~Gesztesy, R.~H{\o}egh-Krohn, H.~Holden,
\emph{Solvable Models in Quantum Mechanics}, Springer, 1988.
\bibitem{Weil}
A.~Weil,
\emph{Sur les formules explicites de la th\'eorie des nombres},
Comm. Sem. Math. Univ. Lund, tome suppl. d\'edi\'e \`a Marcel Riesz (1952), 252--265.
\end{thebibliography}

\end{document}